\chapter{职业有害因素的监督与控制}
\setcounter{chapter}{11}
\chapterintro{
\textbf{【教学大纲要求】}

\imp{掌握：}职业卫生标准（意义与内容、车间空气中有害物质接触限值——MAC、TLV-TWA、PC-STEL；制订依据；生物接触限值；化学致癌物职业接触限值；职业卫生标准的应用）；工业通风（类型和原理）；个人防护用品（不同防护用品的特点、使用注意事项）；工业通风系统的测定与评价。

\imp{熟悉：}职业卫生法规与监督；作业场所采光与照明；作业场所健康促进。
}

\section{职业卫生标准}

\subsection{职业卫生标准的分类与应用}

\begin{defbox}
\textbf{制定职业卫生标准的原则：}
\begin{enumerate}[leftmargin=*]
    \item 保障劳动者不发生急性及慢性职业性损害
    \item 选用最敏感的指标（最敏感人群、最敏感观察指标）
    \item 经济技术上可行
\end{enumerate}
\end{defbox}

\subsection{职业接触限值的应用}

\begin{itemize}[leftmargin=*]
    \item MAC：适用于急性毒性较大的物质，评价瞬时暴露
    \item PC-TWA：适用于慢性毒性物质，评价长期平均暴露水平
    \item PC-STEL：与PC-TWA配合使用，控制短期高浓度暴露
    \item 化学致癌物的职业接触限值：对遗传毒性致癌物，力争将接触水平降至最低
\end{itemize}

\section{工业通风}

\begin{defbox}
\textbf{工业通风}：通过合理组织车间内外的气流，控制生产过程中产生的有害气体、蒸气、粉尘等，使其浓度降低到国家卫生标准以下的工程技术措施。

\textbf{通风类型：}
\begin{enumerate}[leftmargin=*]
    \item \imp{自然通风}：利用风压和热压进行通风（经济节能）
    \item \imp{机械通风}：利用通风机进行通风
    \begin{itemize}
        \item 全面通风：对整个车间的空气进行交换
        \item 局部通风（优先采用）：局部排风（在有害物质发生源就地将其排出）和局部送风
    \end{itemize}
\end{enumerate}

\textbf{局部排风系统组成}：吸尘（气）罩→风管→净化设备→风机→排气筒
\end{defbox}

\section{个人防护用品}

\begin{defbox}
\textbf{个人防护用品（PPE）}：劳动者在生产过程中为防御物理、化学、生物等有害因素伤害人体而穿戴和使用的各种物品。

\textbf{分类：}
\begin{enumerate}[leftmargin=*]
    \item \imp{呼吸防护器}：防尘口罩、防毒面具、供气式呼吸器
    \item \imp{防护服}：防尘服、防毒服、防热服、防辐射服
    \item \imp{防护手套}：防振手套、防酸碱手套、绝缘手套
    \item \imp{防护鞋}：防砸鞋、绝缘鞋、耐酸碱鞋
    \item \imp{护耳器}：耳塞、耳罩
    \item \imp{护目镜和面罩}：防冲击、防化学物飞溅、防辐射
\end{enumerate}

\textbf{使用注意事项：}
\begin{itemize}[leftmargin=*]
    \item 根据危害因素类型选择合适的防护用品
    \item 正确佩戴和使用
    \item 定期检查和更换
    \item 加强维护和保养
\end{itemize}
\end{defbox}

\section{作业环境管理}

\begin{itemize}[leftmargin=*]
    \item 作业场所采光照明应符合标准
    \item 设置应急救援设施和警示标识
    \item 定期检测作业场所有害因素浓度
    \item 建立职业卫生档案
\end{itemize}

% ----- 复习题 -----
\section{复习题}

\begin{enumerate}[leftmargin=*, itemsep=8pt]
    \item \textbf{【名词解释】}工业通风；局部排风；个人防护用品（PPE）
    \item \textbf{【简答题】}简述工业通风的类型及适用条件。
    \item \textbf{【简答题】}简述个人防护用品的分类和使用注意事项。
    \item \textbf{【简答题】}简述职业卫生标准的制定原则。
\end{enumerate}

\subsection*{参考答案要点}

\begin{enumerate}[leftmargin=*]
    \item \textbf{工业通风}：通过合理组织气流控制有害气体的工程技术措施。\textbf{局部排风}：在有害物发生源就地将其排出的通风方式。\textbf{PPE}：防御有害因素伤害人体而穿戴使用的各种物品。
    \item 自然通风（利用风压和热压，经济）和机械通风（全面通风和局部通风）。有害物质源固定时优先采用局部排风。
    \item 分类：呼吸防护器、防护服、防护手套、防护鞋、护耳器、护目镜和面罩。注意：按危害因素选型、正确佩戴、定期检查更换、加强维护保养。
    \item 三项原则：保障不发生急慢性职业损害；选用最敏感指标（最敏感人群、最敏感观察指标）；经济技术上可行。
\end{enumerate}

\newpage

% =====================================================================
%  CHAPTER 12: 主要行业的职业卫生
% =====================================================================
\newpage